% !PRINTED 568
\weberpage{568}
%<<CONT
the sum of two squares of real linear forms. The product
$\eta \eta'$ then appears with the factor 2 in the function $f$.
Hence $f$ is in all circumstances a positive form.

But by (7), (8), $f$ can also be regarded as a quadratic form
of the variables $\eta_1, \eta_2, \ldots, \eta_n$,
\weq{9}{f = F(\eta_1, \eta_2, \ldots, \eta_n) = \sum^{i,k} a_{i,k}\,\eta_i \eta_k}
and in it, if $\eta_i$ is real, $a_{i,i} = 1$, and if
$\eta_i, \eta_k$ is a conjugate imaginary pair, $a_{i,k} = 1$. All the remaining
coefficients are $= 0$. The determinant of $F$ is therefore $= \pm 1$,
because in every row and in every column only one element
different from zero, with the value 1, occurs. By Bd. I,
\S.\,56 the determinant of the form $f$ is equal to the product of the
determinant of $F$ and the square of the determinant of the
linear forms $\eta$; hence by (5)
\weq{10}{= \pm [N_a(\varphi)]^2 \Delta,}
where the sign is to be chosen so that the value comes out positive.

By \S.\,155, (24) one can now put for the $x_i$ such integral
rational numbers that
\weq{11}{f(x_1, x_2, \ldots, x_n) < \Theta_n\, n \sqrt[n]{\pm \Delta N_a(\varphi)^2}}
holds, where $\Theta_n$, if $n$ is greater than 1, is a proper fraction and
always less than $\sqrt[n]{0{,}404}$. If one puts these values of $x_i$ into
(4), then $\eta$ passes over into an integer $\beta$ divisible by $\varphi$,
and with this we therefore have the theorem:

2. \gesperrt{If $\varphi$ is any integral functional in $\Omega$, then there}
\gesperrt{is an integer $\beta$ in $\Omega$ divisible by $\varphi$ with}
\gesperrt{the property that}
\weq{12}{S\,|\eta|^2 < \Theta_n\, n \sqrt[n]{\pm \Delta [N_a(\varphi)]^2}}
\gesperrt{holds, where $\Delta$ is the fundamental number of the field and}
\gesperrt{$\pm \Delta$ is positive.}

The product of the conjugate values $\eta$ is equal to the norm
of $\eta$, and hence by Lemma 1. [with regard to
(7) and (8)]
\weq{13}{n \sqrt[n]{[N(\eta)]^2} \leqq S\,|\eta|^2,}
and so by (12) there follows the inequality relating to the
absolute values:
\weq{14}{N(\eta)^2 < \Theta_n^n\, [N_a(\varphi)]^2\, \Delta.}
Since now $\eta$ is divisible by $\varphi$, we can
\[ \varepsilon \eta = \varphi \psi, \quad \pm N(\eta) = N_a(\varphi)\, N_a(\psi) \]
%CONT>>
